\documentclass[11pt]{article}
\usepackage[margin=1in]{geometry}
\usepackage{hyperref}
\usepackage{enumitem}
\usepackage{amsmath,amssymb,amsthm,mathtools}
\usepackage[T1]{fontenc}
\usepackage[utf8]{inputenc}
\title{\textbf{Theory-mode report: Edge Universality for Inhomogeneous Random Matrices II: Markov Chain Comparison and Critical Statistics}}
\author{Generated from Scholar Inbox digest}
\date{Digest date: 2026-04-23}
\begin{document}
\maketitle
\section*{Paper information}
\begin{itemize}[leftmargin=1.5em]
\item \textbf{Title:} Edge Universality for Inhomogeneous Random Matrices II: Markov Chain Comparison and Critical Statistics
\item \textbf{Authors:} Dang-Zheng Liu, Guangyi Zou
\item \textbf{arXiv:} \href{https://arxiv.org/abs/2604.20215}{2604.20215}
\item \textbf{Scholar digest score:} 95.0
\end{itemize}
\section*{Executive summary}
This report summarizes the highest-scoring paper in the Scholar Inbox digest dated \textbf{2026-04-23}. The abstract states the core claim as follows: \emph{The first paper in this series introduced a \textbackslash\{\}emph\{short-to-long mixing\} condition that captures mean-field GOE/GUE edge universality in the supercritical sparsity regime, for symmetric/Hermitian random matrices with independent entries and a Markov variance profile. This condition reduces the universality problem to the mixing properties of the underlying Markov chains.
  In this paper, we develop new \textbackslash\{\}emph\{short-to-long comparison\} conditions that extend the analysis to the subcritical and critical sparsity regimes. Specifically, we prove that two inhomogeneous random matrices exhibit the same universal edge statistics whenever their variance-profile Markov chains are comparable, regardless of the fine details of the matrix entries. To illustrate the power of our Markov chain comparison theorem, we derive the spectral edge statistics for several prototypical models: random band matrices}.

\section*{Problem setup and intuition}
The introduction reinforces this lens: the paper appears motivated by a gap between practical behavior and theoretical understanding.

\section*{Proof architecture}
The technical core appears to build the surrogate quantity and the chain of lemmas needed to propagate local control into the final theorem: the paper follows a standard theory-paper structure with structural reduction plus quantitative bounds.

\section*{Takeaway}
The key reusable lesson is to define the right object, find the structural reduction, and quantify the error terms clearly.
\end{document}
